\documentclass[a4paper,11pt]{article}

%%% Packages
\usepackage[utf8]{inputenc}
\usepackage[T1]{fontenc}
\usepackage{amsmath,amssymb,amsthm}
\usepackage{mathtools}
\usepackage{booktabs}
\usepackage{array}
\usepackage{enumitem}
\usepackage{hyperref}
\usepackage[margin=2.5cm]{geometry}
\usepackage{xcolor}

%%% Theorem environments
\newtheorem{theorem}{Theorem}[section]
\newtheorem{lemma}[theorem]{Lemma}
\newtheorem{proposition}[theorem]{Proposition}
\newtheorem{corollary}[theorem]{Corollary}
\newtheorem{claim}[theorem]{Claim}
\theoremstyle{definition}
\newtheorem{definition}[theorem]{Definition}
\newtheorem{example}[theorem]{Example}
\newtheorem{remark}[theorem]{Remark}

%%% Notation macros
\newcommand{\ALC}{\ensuremath{\mathcal{ALC}}}
\newcommand{\ALCI}{\ensuremath{\mathcal{ALCI}}}
\newcommand{\ALCIRCC}[1]{\ensuremath{\mathcal{ALCI}_{\mathit{RCC#1}}}}
\newcommand{\DR}{\ensuremath{\mathsf{DR}}}
\newcommand{\PO}{\ensuremath{\mathsf{PO}}}
\newcommand{\EQ}{\ensuremath{\mathsf{EQ}}}
\newcommand{\PP}{\ensuremath{\mathsf{PP}}}
\newcommand{\PPI}{\ensuremath{\mathsf{PPI}}}
\newcommand{\inv}{\ensuremath{\mathsf{inv}}}
\newcommand{\comp}{\ensuremath{\mathsf{comp}}}
\newcommand{\cl}{\ensuremath{\mathsf{cl}}}
\newcommand{\Tp}{\ensuremath{\mathsf{Tp}}}
\newcommand{\Tri}{\ensuremath{\mathsf{Tri}}}
\newcommand{\Prof}{\ensuremath{\mathsf{Prof}}}
\newcommand{\Col}{\ensuremath{\mathsf{Col}}}
\newcommand{\Sig}{\ensuremath{\mathsf{Sig}}}
\newcommand{\Req}{\ensuremath{\mathsf{Req}}}
\newcommand{\Pay}{\ensuremath{\mathsf{Pay}}}
\newcommand{\old}{\ensuremath{\mathsf{old}}}
\newcommand{\Roles}{\ensuremath{\mathcal{R}}}
\newcommand{\EXPTIME}{\textsc{ExpTime}}
\newcommand{\PSPACE}{\textsc{PSpace}}
\newcommand{\NP}{\textsc{NP}}

\title{Assessment of the Profile-Cached Blocking Calculus\\
  for \texorpdfstring{$\ALCIRCC{5}$}{ALCI\_RCC5} under Strong EQ\\[0.4em]
  \large Response to GPT's draft series (April 2, 2026)}

\author{Claude\thanks{Claude is an AI assistant by Anthropic.
  This assessment was prepared at the request of Michael Wessel
  (\texttt{miacwess@gmail.com}).  It evaluates four draft notes
  by GPT proposing a profile-cached blocking calculus for
  $\ALCIRCC{5}$ under strong EQ semantics.}}

\date{April 2026}

\begin{document}
\maketitle

\begin{center}
\fbox{\parbox{0.92\textwidth}{%
\textbf{Disclaimer.}\; This assessment was authored entirely by Claude
(Anthropic), an AI assistant.  The evaluation and the claims made
herein have \textbf{not been verified by human domain experts}.  Both
the assessed drafts and this response are discussion pieces.}}
\end{center}

\bigskip

\begin{abstract}
GPT has produced a series of four draft notes proposing a
profile-cached global blocking calculus for $\ALCIRCC{5}$ under strong
EQ semantics.  The series exhibits an impressive self-correction cycle:
an initial conditional result is refined through three revisions that
correct a false lemma, provide an explicit finite signature system,
and ultimately claim a full replay theorem based on a meet-semilattice
structure on RCC5 witness labels.

We assess each note in turn.  The \textbf{local combinatorial results
are correct and novel}: the counterexample to flat normalization, the
coherent-block normalization lemma, the meet-polymorphism property of
the RCC5 composition table, and the robust colorwise normalization
theorem all withstand scrutiny and computational verification.
However, the \textbf{global unraveling argument in the final note has
a genuine gap}: the unraveling changes the predecessor color structure
(multiplicities, inter-color relations, and even the color alphabet
itself), so the replay theorem---proved for nodes with equal
signatures in the finite tableau---cannot be directly applied to the
unraveling's copies.

We exhibit this gap concretely on the canonical infinite-model concept
$C_\infty = \exists\PP.\top \sqcap \forall\PP.\exists\PP.\top$.  We
then discuss how the gap might be closed, noting that a
patchwork-based argument for cross-chain edge assignment could
complement GPT's local machinery.
\end{abstract}

\tableofcontents

%% ================================================================
\section{Introduction}
\label{sec:intro}

The decidability of concept satisfiability in $\ALCIRCC{5}$ under
strong EQ semantics has been an open problem since Wessel's
report~\cite{Wessel2003}.  Multiple approaches have been attempted:
\begin{itemize}[nosep]
\item A quasimodel method with type elimination~\cite{CompanionPaper},
  which has the \emph{extension gap} in the soundness direction (open
  quasimodel need not yield a model);
\item A contextual tableau~\cite{GPTContextual}, whose completeness
  conjecture FW$(C,N)$ is false~\cite{FWCounterexample};
\item A direct model construction from open completion
  graphs~\cite{ExtensionGap}, which identifies the root cause as the
  self-absorption failure $\PPI \notin \comp(\DR, \PPI) = \{\DR\}$
  and proposes a resolution (with caveats on the path-consistency
  argument).
\end{itemize}

GPT now proposes a fourth approach: a \textbf{profile-cached global
blocking calculus} in the style of anywhere blocking / global caching,
adapted to complete-graph semantics.  The approach is developed across
four draft notes:
\begin{enumerate}[nosep]
\item \textbf{Blocking Draft} --- initial proposal, conditional on a
  classwise witness-normalization lemma;
\item \textbf{Blocking Revised} --- self-correction: flat
  normalization is false; repaired via coherent predecessor blocks;
\item \textbf{Explicit Signatures} --- concrete depth-indexed
  signature system with finite-index proof;
\item \textbf{Replay Final} --- claims full replay theorem via
  meet-semilattice normalization, yielding decidability.
\end{enumerate}

This note assesses the four papers in order, verifies the local
algebraic claims computationally, and identifies the remaining gap.


%% ================================================================
\section{Paper 1: The Initial Draft}
\label{sec:paper1}

\subsection{Architecture}

The draft proposes a global blocking scheme for complete-graph
tableaux.  Each saturated node~$x$ receives a \emph{contextual
profile}
\[
  \Prof(x) = \bigl(\lambda(x),\; \mathsf{Ctx}(x)\bigr),
\]
where $\lambda(x)$ is the saturated Hintikka type and
$\mathsf{Ctx}(x)$ is a finite quotient of the older neighborhood,
collapsing nodes with identical \emph{context signatures}
$(\lambda(u), \rho(u,x), \Gamma_{\rho(u,x)}(\lambda(u)))$.

A node $y$ is blocked by an older node $x$ when $\Prof(y) =
\Prof(x)$.  Existential witnesses are recorded as \emph{witness
recipes}: a target type plus a classwise assignment of RCC5 labels
from each context class to the new witness.  Termination follows from
the finiteness of the profile universe.

\subsection{Assessment}

The architecture is sound and well-motivated.  The observation that
blocking must be global (not ancestor-based) is correct: in
complete-graph models, every older node constrains every newer node,
so tree-based blocking is inapplicable.

The paper is honest about its status: all results are conditional on
a \textbf{classwise witness-normalization lemma} (Lemma~4.1 in the
draft), which states that witnesses can be chosen so that their
behavior depends only on context classes, not individual
representatives.  This is correctly identified as the key proof
obligation.

\textbf{Verdict: clean conditional result.}  The metatheory is
correct modulo the stated assumption.


%% ================================================================
\section{Paper 2: The Self-Correction}
\label{sec:paper2}

\subsection{The counterexample to flat normalization}

The most important contribution of Paper~2 is
Proposition~4.2: flat classwise normalization is \textbf{false}.

\begin{example}[GPT's counterexample, verified]
\label{ex:flat-false}
Node $x$ carries $\exists\DR.(p \sqcap q \sqcap r)$.  Older nodes
$a_1, a_2, b_1, b_2, c$ are all $\DR$-related to~$x$.  Types:
\[
  \tau_A\!: \forall\PO.\neg p,\, \forall\PP.\neg p,\,
            \forall\PPI.\neg p;
  \quad
  \tau_B\!: \forall\PO.\neg q,\, \forall\PP.\neg q;
  \quad
  \tau_C\!: \forall\DR.\neg r,\, \forall\PP.\neg r,\,
            \forall\PPI.\neg r.
\]
All inter-predecessor edges are $\DR$ except $b_2\;\PPI\;c$.
Target type $t = \{p, q, r, \top\}$, witness labels $\eta(a_i) =
\DR$, $\eta(b_1) = \DR$, $\eta(b_2) = \PPI$, $\eta(c) = \PO$.

\smallskip\noindent
\textbf{Verification.}  The payload constraints force:
\begin{itemize}[nosep]
\item From $\tau_A$: only $\DR$ is allowed (all other payloads
  contain $\neg p$, but $p \in t$).
\item From $\tau_B$: only $\DR$ or $\PPI$ are allowed ($\PO$ and
  $\PP$ payloads contain $\neg q$).
\item From $\tau_C$: only $\PO$ is allowed ($\DR$, $\PP$, $\PPI$
  payloads contain $\neg r$).
\end{itemize}
The triangle constraint for $b_1$ (with $c$): $\Tri(\DR, \PO,
\eta(b_1))$ requires $\eta(b_1) \in \comp(\DR, \PO) = \{\DR, \PO,
\PP\}$.  Combined with the payload restriction to $\{\DR, \PPI\}$,
we get $\eta(b_1) = \DR$ (forced).

The triangle constraint for $b_2$ (with $c$): $\Tri(\PPI, \PO,
\eta(b_2))$ requires $\eta(b_2) \in \comp(\PPI, \PO) = \{\PO,
\PPI\}$.  Combined with the payload restriction to $\{\DR, \PPI\}$,
we get $\eta(b_2) = \PPI$ (forced).

Since $b_1$ and $b_2$ share the same flat class but are forced to
different labels, no flat-classwise normalization exists.  $\square$
\end{example}

\begin{remark}[Connection to our self-absorption analysis]
\label{rem:self-absorption}
The root cause is precisely the self-absorption failure we identified
independently: $\PPI \notin \comp(\DR, \PPI) = \{\DR\}$.  The edge
$b_2\;\PPI\;c$ (versus $b_1\;\DR\;c$) changes the triangle
constraint, forcing different witness labels despite identical flat
classes.  GPT's counterexample and our algebraic analysis converge on
the same structural rigidity in the RCC5 composition table.
\end{remark}


\subsection{Coherent-block normalization}

Paper~2 introduces \textbf{coherent predecessor blocks}: a refinement
of flat classes where inter-block relations are uniform and
intra-block relations are uniform symmetric ($\DR$ or $\PO$).

\begin{proposition}[Coherent-block normalization --- verified]
Every locally admissible symbolic extension $(t, \eta)$ can be
normalized to $(t, \eta')$ that is constant on each coherent block,
using only labels already realized in the original extension.
\end{proposition}

The proof proceeds by cases on the intra-block label:
\begin{itemize}[nosep]
\item \textbf{$\PO$-blocks}: Every label $r$ is self-stable under
  $\PO$: $\Tri(\PO, r, r)$ holds for all $r \in \Roles$.  (Verified:
  $r \in \comp(\PO, r)$ for all $r$.)  Any realized label works.

\item \textbf{$\DR$-blocks}: $\PPI$ is the only unstable label:
  $\Tri(\DR, \PPI, \PPI)$ fails because $\PPI \notin \comp(\DR, \PPI)
  = \{\DR\}$.  So in a non-singleton $\DR$-block, not all members can
  carry $\PPI$ (two $\PPI$-labeled members would violate the triangle
  condition).  Therefore some label in $\{\DR, \PO, \PP\}$ is
  realized, and all three are $\DR$-stable.

\item \textbf{Inter-block}: By coherence (uniform edge $q_{BC}$),
  one can pick representatives $u \in B$, $v \in C$ realizing the
  chosen labels and appeal to the original admissibility.
\end{itemize}

\textbf{Verdict: correct and clean.}  This is the right local repair.


%% ================================================================
\section{Paper 3: Explicit Signatures and Finite Index}
\label{sec:paper3}

Paper~3 (the ``explicit signatures'' addendum) makes the depth-indexed
signature construction concrete and proves the finite-index lemma.
The construction uses:
\begin{itemize}[nosep]
\item Rank-0 signatures = propositional types $\mathsf{tp}_0(x)$.
\item Rank-$(k{+}1)$ predecessor colors built from rank-$k$
  signatures, incoming relation, truncated payload, refined through
  $k{+}1$ rounds of relation-set profiling.
\item Context tables: multiplicity map $\mu$ (values in $\{0, 1,
  +\}$) and relation-set table~$Q$.
\item Projected witness footprints recording which labels occur on
  coherent blocks of each final color.
\end{itemize}

The finite-index lemma is correct: at each rank, all components range
over finite domains, so the signature space is finite.

The paper is deliberately cautious: it proves finiteness of the
signature system but does \textbf{not} claim the replay theorem.  It
explicitly states that turning projected footprints into functional
replay recipes is the remaining step.

\textbf{Verdict: correct.}  The finite-index part of the blocking
program is settled.


%% ================================================================
\section{Paper 4: The Meet-Semilattice and Replay Theorem}
\label{sec:paper4}

Paper~4 claims to close the last gap.  The key innovations are:
\begin{enumerate}[nosep]
\item Predecessor colors now include the \textbf{full payload vector}
  $(\Gamma_R(\lambda(u)) \cap \cl_k(C))_{R \in \Roles}$, not just
  the incoming-edge payload.
\item Witness footprints are sharpened to \textbf{robust colorwise
  recipes} $\chi : A_{k+1} \to \Roles$ using a meet-normalization.
\end{enumerate}

\subsection{The meet-semilattice structure (verified)}
\label{sec:meet}

The four non-EQ RCC5 labels are partially ordered by the diamond:
\[
  \DR < \PO < \PP, \qquad \DR < \PO < \PPI,
  \qquad \PP \text{ and } \PPI \text{ incomparable}.
\]
The meet (greatest lower bound) gives:
\[
  \DR \wedge r = \DR, \quad
  \PO \wedge \PP = \PO, \quad
  \PO \wedge \PPI = \PO, \quad
  \PP \wedge \PPI = \PO.
\]

\begin{lemma}[Meet polymorphism --- computationally verified]
\label{lem:meet-poly}
For every $q \in \Roles$, the relation
$R_q = \{(r,s) \mid \Tri(q, s, r)\}$
is closed under componentwise meet.
\end{lemma}

\noindent
We verified all $4 \times 16 \times 16 = 1{,}024$ cases
computationally against the full RCC5 composition table under strong
EQ.  Zero violations.

\begin{lemma}[No $\PP/\PPI$ within one final color --- correct]
If two predecessors of the same final color were related by~$\PP$,
every node of that color would need an outgoing $\PP$-edge within the
color class.  But $\PP$ on a finite set is a strict partial order,
which has maximal elements with no outgoing $\PP$-edge.
Contradiction.
\end{lemma}

\begin{lemma}[Realized labels cannot mix $\PP$ and $\PPI$ --- correct]
If coherent blocks of the same color carried both $\PP$ and $\PPI$ as
witness labels, the triangle condition would require
$\Tri(q, \PPI, \PP)$ for $q \in \{\DR, \PO\}$ (the only possible
intra-color relations by the previous lemma).  Verified: $\PP \notin
\comp(q, \PPI)$ for both $q = \DR$ and $q = \PO$.  Contradiction.
\end{lemma}

\begin{theorem}[Robust colorwise normalization --- correct]
For each final color~$a$, the meet
$m(a) = \bigwedge \{\eta(B) \mid B \text{ coherent block of
color } a\}$
is in the realized set (by the previous lemma) and yields a locally
admissible extension.  The proof uses the covering argument: for each
pair of colors $a, b$ and relation $q \in Q(a,b)$, choose a covering
family of block pairs, apply admissibility to each, and take the
componentwise meet (which stays in $R_q$ by meet polymorphism).
\end{theorem}

\textbf{Verdict on local results: correct, elegant, and novel.}  The
meet-semilattice structure is a genuinely nice algebraic insight that
we had not identified in our own analysis.

\subsection{A missing admissibility constraint (minor)}
\label{sec:missing-constraint}

The definition of \emph{locally admissible symbolic extension} in
all four papers checks only:
\begin{enumerate}[label=(\roman*),nosep]
\item Payload preservation: $\Gamma_{\eta(u)}(\lambda(u)) \subseteq
  t$ for $u \in \old(x)$.
\item Pairwise triangles: $\Tri(\rho(u,v), \eta(v), \eta(u))$ for
  distinct $u, v \in \old(x)$.
\end{enumerate}

\noindent
However, when a witness~$y$ is created for $\exists S.E$ at node~$x$,
there is also the triangle $(u, x, y)$ for each $u \in \old(x)$:
\begin{equation}
  \label{eq:missing-tri}
  \eta(u) \in \comp(\rho(u,x),\; S).
\end{equation}
This constraint is not listed in the admissibility definition.

\medskip
This is a \textbf{presentation issue, not a correctness issue}.  We
verified computationally that for every $q, S \in \Roles$, the set
$\comp(q, S)$ is closed under the diamond meet.  Therefore the
meet-normalization preserves this constraint even though it is not
formally tracked.  Moreover, model-extracted recipes automatically
satisfy~\eqref{eq:missing-tri}.  The fix is simply to add this
constraint to the definition.


%% ================================================================
\section{The Unraveling Gap}
\label{sec:gap}

The critical issue is Theorem~6.2 of Paper~4 (the ``Blocked
Unraveling Theorem''), which claims that an open saturated blocked
branch can be unraveled into a genuine model.  The proof is a
one-paragraph sketch asserting that ``standard induction on formula
depth'' suffices.

The replay theorem (Theorem~6.1) is correct: if two nodes in the
finite tableau have equal rank-$d$ signatures, then robust recipes at
one can be applied at the other.  The gap is that \textbf{nodes in the
unraveling need not have the same signature as the representative they
are copies of}.

\subsection{Concrete demonstration: $C_\infty$}

Consider $C_\infty = \exists\PP.\top \sqcap
\forall\PP.\exists\PP.\top$, which requires an infinite model (every
element must have a $\PP$-successor, and so must every
$\PP$-descendant).

\paragraph{The finite tableau.}
\begin{itemize}[nosep]
\item Node $r_0$: type $\tau = \{\exists\PP.\top,\;
  \forall\PP.\exists\PP.\top\}$, no predecessors ($\old(r_0) =
  \emptyset$).
\item Node $r_1$: $\PP$-witness for $r_0$, type~$\tau$ (forced by the
  $\forall\PP$-propagation), predecessor set $\old(r_1) = \{r_0\}$.
\item $r_1$ is blocked by $r_0$ (same profile in a suitable sense).
\end{itemize}
The recipe at $r_1$ for the existential $\exists\PP.\top$ assigns a
label to its sole predecessor~$r_0$.  The triangle $(r_0, r_1,
\text{witness})$ with $\rho(r_0, r_1) = \PP$ and $S = \PP$ forces
$\eta(r_0) \in \comp(\PP, \PP) = \{\PP\}$, so $\eta(r_0) = \PP$.

\paragraph{The unraveling.}
The unraveling creates an infinite $\PP$-chain:
\[
  c_0 \xrightarrow{\PP} c_1 \xrightarrow{\PP} c_2
  \xrightarrow{\PP} c_3 \xrightarrow{\PP} \cdots
\]
At step~2, $c_2$ is created as a $\PP$-witness for $c_1$.  The recipe
assigns $\rho(c_0, c_2) = \PP$ (from the recipe at $r_1$).

At step~3, $c_3$ is created for $c_2$.  Now $\old(c_2) = \{c_0,
c_1\}$.  Both predecessors have initial color
$(\tau, \PP, \Pay(\tau))$ relative to~$c_2$ (since $\rho(c_0, c_2) =
\PP$ and $\rho(c_1, c_2) = \PP$).  After one refinement round:
\begin{itemize}[nosep]
\item $c_0$ sees $c_1$ (same initial color) via $\PP$, so its
  relation-set profile includes $\PP$ to its own initial color.
\item $c_1$ sees $c_0$ (same initial color) via $\PPI$, so its
  relation-set profile includes $\PPI$ to the same initial color.
\end{itemize}
Since $\PP \neq \PPI$, the two predecessors get \textbf{different
final colors} after one refinement round.

\paragraph{The mismatch.}
In the finite tableau, the representative $r_1$ had \emph{one}
predecessor ($r_0$) of one color, with multiplicity~1.  In the
unraveling, $c_2$ has \emph{two} predecessors of \emph{two different}
colors, each with multiplicity~1.  The multiplicity maps differ; the
relation-set tables differ (the unraveling has a $\PP$/$\PPI$
inter-color relation that did not exist in the finite tableau); and
the recipe at~$r_1$ has no assignment for the new color because that
color did not exist in~$r_1$'s context.

At step~$n$, the situation is worse: $c_n$ has $n$ predecessors, which
after color refinement split into $O(n)$ distinct final colors.  The
recipe, designed for a bounded neighborhood, cannot cover them.

\paragraph{Conclusion.}
The replay theorem applies to nodes with \emph{equal signatures}.  But
copies in the unraveling accumulate predecessors whose color structure
diverges from the finite tableau's.  The replay theorem's
prerequisites are not met, and the unraveling construction does not
produce a certified model.

\subsection{Why this gap is non-trivial}
\label{sec:why-hard}

The gap is structural, not cosmetic.  In tree-based DL tableaux,
unraveling is straightforward because each node's neighborhood is its
parent and children---a bounded, predictable structure.  In
complete-graph tableaux for $\ALCIRCC{5}$, every node is connected to
every other node.  The unraveling must assign RCC5 labels to
\emph{all pairs} of copies, including cross-chain pairs whose
existence was not anticipated by the finite tableau's recipes.

This is the same fundamental difficulty that affects all three
previous approaches:
\begin{itemize}[nosep]
\item In the quasimodel approach, it manifests as the extension gap
  (abstract quasimodels need not yield models).
\item In the contextual tableau, it manifests as the failure of
  FW$(C,N)$ (local states of bounded width cannot capture
  the global model).
\item In our direct model construction, it manifests as the
  path-consistency argument for the cross-subtree constraint network
  (Cases~B and~C).
\end{itemize}

GPT's profile-cached approach elegantly handles the \emph{local}
problem (what labels to assign from predecessors to a new witness)
but does not fully address the \emph{global} problem (what labels to
assign between copies that were not simultaneously present in the
finite tableau).


%% ================================================================
\section{Possible Repairs}
\label{sec:repairs}

The gap is specific and the local machinery is strong.  Several
repair strategies are conceivable.

\subsection{Patchwork-based cross-chain assignment}
\label{sec:repair-patchwork}

Our companion paper~\cite{ExtensionGap} handles cross-subtree edges
by setting up a \emph{disjunctive RCC5 constraint network} over
all pairs of copies and appealing to \textbf{full RCC5 tractability}
(Renz \& Nebel~\cite{RenzNebel1999}): a path-consistent disjunctive
RCC5 network is globally consistent.

A hybrid approach could work as follows:
\begin{enumerate}[nosep]
\item Use GPT's profile-cached blocking for \textbf{termination} and
  \textbf{local recipe generation}.
\item When unraveling, assign \textbf{tree edges} (parent-to-child
  via recipes) as in GPT's construction.
\item For \textbf{cross-chain edges} (between copies not in a
  parent-child relationship), set up a disjunctive constraint network
  whose constraints come from the composition table applied along the
  tree path connecting the two copies.
\item Prove \textbf{path-consistency} of this network using the
  tableau's properties and the meet-normalization results.
\item Apply full RCC5 tractability to obtain a globally consistent
  atomic refinement.
\end{enumerate}

This would marry GPT's local elegance with the global consistency
guarantee from the patchwork property.

\subsection{Weakened replay with monotone signatures}
\label{sec:repair-monotone}

Alternatively, one could define a notion of \emph{signature
subsumption} such that a copy~$c$ in the unraveling has a signature
that \emph{subsumes} its representative's signature (more colors,
richer context).  If one can show that recipes designed for the
simpler signature remain admissible under the richer one---using the
monotonicity of the meet-normalization and the
$\comp(q,S)$-meet-closure we verified---then the replay theorem could
be extended to subsumption rather than equality.

This seems plausible but would require careful formalization of the
subsumption relation and its interaction with the recipe admissibility
conditions.

\subsection{Truncated unraveling with depth-bounded predecessors}

A third option: instead of unraveling the full cache, build a
\emph{truncated} unraveling where each copy only remembers
predecessors up to some bounded depth in the unraveling tree.  The
truncation depth would be chosen so that the signature stays within
the finite alphabet.  Cross-chain edges beyond the truncation horizon
would be assigned by a default rule (e.g., $\DR$, which is always
composition-consistent with itself).  The correctness of the truth
lemma under truncation would need to be verified.


%% ================================================================
\section{Summary of Verification Results}
\label{sec:summary}

\begin{center}
\renewcommand{\arraystretch}{1.25}
\begin{tabular}{@{}lcc@{}}
\toprule
\textbf{Claim} & \textbf{Paper} & \textbf{Status} \\
\midrule
Flat normalization is false (counterexample) & 2 & \textbf{Correct} \\
Coherent-block normalization & 2 & \textbf{Correct} \\
Explicit finite-index signature system & 3 & \textbf{Correct} \\
Meet polymorphism ($R_q$ closed under $\wedge$) & 4 &
  \textbf{Verified} (1,024 cases) \\
No $\PP/\PPI$ within one final color & 4 & \textbf{Correct} \\
Realized labels don't mix $\PP$ and $\PPI$ & 4 & \textbf{Correct} \\
Robust colorwise normalization & 4 & \textbf{Correct} \\
$\comp(q,S)$ is meet-closed (missing constraint) & --- &
  \textbf{Verified} (64 cases) \\
Replay theorem (finite tableau) & 4 & \textbf{Correct} \\
\midrule
Blocked unraveling theorem & 4 & \textbf{Gap} \\
\bottomrule
\end{tabular}
\end{center}


%% ================================================================
\section{Comparison of Approaches}
\label{sec:comparison}

\begin{center}
\renewcommand{\arraystretch}{1.25}
\begin{tabular}{@{}p{4.2cm}p{4.8cm}p{4.8cm}@{}}
\toprule
& \textbf{GPT: profile-cached blocking}
& \textbf{Claude: direct model construction} \\
\midrule
Local arguments
& Meet normalization, coherent blocks --- \textbf{proved}
& Self-absorption, forced $\DR$ edges --- \textbf{proved} \\[3pt]
Global model construction
& Unraveling from finite cache --- \textbf{gap}
& Tree unraveling + constraint network --- \textbf{gap} (Cases B/C) \\[3pt]
Key algebraic insight
& Meet-semilattice on RCC5 labels
& Self-absorption failure $\comp(\DR,\PPI) = \{\DR\}$ \\[3pt]
Use of patchwork / tractability
& Not used
& Central (full RCC5 tractability) \\[3pt]
Blocking / termination
& Explicit (profile-cached, proved)
& Inherited from original tableau \\
\bottomrule
\end{tabular}
\end{center}

\noindent
The two approaches have \textbf{complementary strengths}.  GPT's meet
normalization elegantly solves the local witness-recipe problem; our
patchwork-based argument addresses the global edge-consistency problem.
A synthesis using GPT's blocking for termination and patchwork for
cross-chain edges in the unraveling could potentially close both
remaining gaps.


%% ================================================================
\section{Conclusion}
\label{sec:conclusion}

GPT's draft series represents strong and largely correct work on a
genuinely difficult problem.  The self-correction cycle---from a
false lemma through coherent blocks to the meet-semilattice---is a
model of honest, iterative proof development.

The local results (coherent-block normalization, meet polymorphism,
robust colorwise normalization, finite-index signatures, replay
theorem for the finite tableau) are all correct, verified, and novel.
The meet-semilattice structure on RCC5 witness labels is an insight
that neither we nor GPT's earlier contextual tableau had identified.

The global unraveling argument, however, has a genuine gap.  The
unraveling changes the predecessor color structure in ways that
prevent direct application of the replay theorem.  This is the same
structural difficulty---in different guise---that affects all
approaches to $\ALCIRCC{5}$ decidability.

We believe the gap is closable, most likely by incorporating a
patchwork-based argument for cross-chain edge assignment.  The
combination of GPT's profile-cached blocking (for termination and
local recipes) with the patchwork property (for global consistency)
appears to be the most promising path toward a complete proof.


%% ================================================================
\begin{thebibliography}{99}

\bibitem{CompanionPaper}
Claude (Anthropic), prompted by M.~Wessel.
\newblock On the decidability of $\mathcal{ALCI}_{\mathit{RCC5}}$ and
  $\mathcal{ALCI}_{\mathit{RCC8}}$: Quasimodels meet the patchwork
  property.
\newblock Discussion paper, April 2026.

\bibitem{ExtensionGap}
Claude (Anthropic), prompted by M.~Wessel.
\newblock Closing the extension gap: Decidability of
  $\mathcal{ALCI}_{\mathit{RCC5}}$ via direct model construction.
\newblock Companion paper, April 2026.

\bibitem{FWCounterexample}
Claude (Anthropic), prompted by M.~Wessel.
\newblock Counterexample to the finite-width extraction conjecture
  FW$(C,N)$ for $\mathcal{ALCI}_{\mathit{RCC5}}$.
\newblock Technical note, March 2026.

\bibitem{GPTContextual}
GPT (OpenAI).
\newblock A contextual tableau calculus for
  $\mathcal{ALCI}_{\mathit{RCC5}}$.
\newblock Draft note, March 2026.

\bibitem{GPTBlocking}
GPT (OpenAI).
\newblock $\mathcal{ALCI}_{\mathit{RCC5}}$ under strong EQ: Replay,
  blocking, and the final meta-theorem.
\newblock Draft note series (4 revisions), April 2, 2026.

\bibitem{RenzNebel1999}
J.~Renz and B.~Nebel.
\newblock On the complexity of qualitative spatial reasoning: A
  maximal tractable fragment of the {Region Connection Calculus}.
\newblock {\em Artificial Intelligence}, 108(1-2):69--123, 1999.

\bibitem{Wessel2003}
M.~Wessel.
\newblock Qualitative spatial reasoning with the
  {$\mathcal{ALCI}_\mathit{RCC}$} family---{F}irst results and
  unanswered questions.
\newblock Technical Report FBI-HH-M-324/03, University of Hamburg,
  2002/2003.

\end{thebibliography}

\end{document}
